\chapter{Fondamenti di probabilità}

\section{Variabili aleatorie}
Una \textbf{variabile aleatoria} (V.A.) è una tripla ($x$, $A_x$, $P_x$), dove il risultato $x$ è il valore di una variabile casuale di un insieme di possibili valori $A_x=\{a_0,a_1,\cdots,a_{n-1}\}$, ognuno con probabilità $P_x=\{p_0, p_1,\cdots,p_{n-1}\}$.
%
Sulla probabilità, valgono i seguenti assiomi:
\begin{enumerate}
	\item \textbf{Non negatività}: Dato un evento $E$ allora $P(E) \ge0$
	\item \textbf{Additività}: Siano $A,B$ 2 V.A. indipendenti, allora $P[A\cup B]=P[A]+P[B]$
	\item \textbf{Normalizzazione}: Sia $E$ l'evento certo, $P[E]=1 $
\end{enumerate}

Da questi assiomi possiamo dimostrare tutti i teoremi sulle probabilità. Possiamo ad esempio dimostrare la probabilità dell'evento impossibile $I$.
$$P[E] = P[E\cup I]=P[E]+P[I]=1+P[I]\Rightarrow P[I] = P[E]-1 = 0$$
Da ora in poi useremo la notazione: $P[A,B] = P[A\cap B] = P[A\wedge B]$

\subsection{Teoremi su Variabili Aleatorie}
Date 2 variabili aleatorie $X,Y$ con supporti $A_x,A_y$ valgono i seguenti teoremi:
\begin{itemize}
	\item \textbf{Teorema della Probabilità Totale}: $$
	P[X=x_i]=\sum_{y\in A_y}P[X=x_i \wedge Y = y_i]$$
	\item \textbf{Probabilità Marginale}: Possiamo ottenere $P(x)$ da $P(x,y)$ tramite una sommatoria: 
	$$
	P(x=a_i) \equiv \sum_{y\in A_Y} P(x=a_i, y) \quad \Rightarrow \quad P(y) = \sum_{x\in A_X} P(x,y)
	$$
	\item \textbf{Probabilità condizionata}: $$P[X|Y]=\frac{P[X\wedge Y]}{P[Y]}$$
	da cui segue: $$P[X=x\wedge Y =y]=P[X=x|Y=y]P[Y=y]$$
	\item \textbf{Teorema di Bayes}: $$P[X|Y] = \frac{P[Y|X]P[X]}{P[Y]}$$
\end{itemize}

\subsubsection{Esempio N.1 - Meccanico}
Alice porta la sua auto dal meccanico. Il meccanico, Bob,
dopo aver effettuato un test computerizzato, sostiene che un pezzo molto costoso debba essere sostituito. Sia \textbf{a} la variabile aleatoria che indica lo stato effettivo del pezzo (a=1 il pezzo è guasto, a=0 il pezzo non è guasto). Si supponga di sapere inoltre che:
\begin{enumerate}
	\item Il test fatto da Bob ha un’affidabilità del 95\%. Nel 95\% dei casi in cui il pezzo è rotto allora la risposta fornita è “il pezzo è rotto”. Nel 95\% dei casi in cui il pezzo non è rotto allora la risposta fornita è “il pezzo non è rotto”;
	\item Solo l’1\% delle vetture (in condizioni analoghe a quella di Alice) ha avuto problemi con il pezzo in questione. 
\end{enumerate}
Qual è la probabilità che il pezzo sia effettivamente guasto?

Indichiamo con $a$ la variabile aleatoria che ci dà lo \textbf{stato effettivo del pezzo}:
\begin{enumerate}
	\item Con $a=1$ il pezzo \textbf{è} guasto, dove $Pr[a=1] = \frac{1}{100}$
	\item Con $a=0$ il pezzo \textbf{non è} guasto, dove $Pr[a=0] = \frac{99}{100}$
\end{enumerate}

e con $b$ la V.A. che ci dà \textbf{l'esito del test}:

\begin{enumerate}
	\item Con $b=1$ il test dice "Rotto".
	\item Con $b=0$ il test dice "OK".
\end{enumerate} 

Sappiamo inoltre che:
\begin{itemize}
	\item $P[b=1|a=1] = \frac{95}{100} \qquad \Rightarrow \qquad P[b=0|a=1] = \frac{5}{100}$
	\item $P[b=0|a=0] = \frac{95}{100} \qquad \Rightarrow \qquad P[b=1|a=0] = \frac{5}{100}$
\end{itemize}

Per trovare la probabilità che il pezzo sia effettivamente guasto dobbiamo trovare:
$$
P[a=1|b=1] =  \frac{P[a=1 \wedge b=1]}{P[b=1]} 
$$

Per prima cosa troviamo $P[a=1 \wedge b=1]$ che sappiamo è uguale a:
$$
P[b=1|a=1]\cdot P[a=1] = \frac{95}{100} \cdot \frac{1}{100}
$$

Successivamente troviamo $P[b=1]$ che sappiamo è uguale a:
$$
P[b=1 \wedge a=1] + P[b=1 \wedge a=0]
$$

Possiamo ricavare $P[b=1 \wedge a=0]$, essendo uguale a:
$$
P[b=1 | a=0] \cdot P[a=0] = \frac{5}{100} \cdot \frac{99}{100}
$$

Riscriviamo il tutto e ricaviamo $P[a=1|b=1]$:
$$
P[a=1|b=1] =  \frac{P[a=1 \wedge b=1]}{P[b=1]} = \frac{\frac{95}{100} \cdot\frac{1}{100}}{\frac{95}{100} \cdot\frac{1}{100} + \frac{5}{100} \cdot \frac{99}{100}} = \frac{95}{95 + 5 \cdot 99} = 0,16
$$

\newpage

\section{Valore Atteso e Varianza}

Il \textbf{Valore atteso} $E[X]$ rappresenta la nostra aspettativa sul valore di una variabile aleatoria. 
\begin{equation}
	E[X] = \sum_{x\in A_x}xP[X=x]
\end{equation}

\textbf{Proprietà del Valore Atteso}:
\begin{itemize}
	\item $E[X,Y] = E[X]+E[Y]$ -- Dimostrazione:\\ $$E[X,Y]=\sum_x\sum_y(x+y)P(x\wedge y) = $$ 
	$$\sum_x \sum_y (xP(x\wedge y))+y(P(x\wedge y)) = $$
	$$ \sum_x x \sum_y P(x,y) + \sum_y y \sum_x P(x,y) = $$
	$$\sum_x x P(x) + \sum_y y P(y) = E[X] +E[Y]$$
	\item $E[cX] = cE[x]$ e $E[X+c] = E[X] + c$
	\item Supponendo $X,Y$ indipendenti $E[XY] = E[X]E[Y]$ -- Dimostrazione:\\
	$$E[XY] = \sum_{x,y} xyP(x\wedge y) = \sum_x\sum_y xy P(X)P(Y)= $$
	$$\sum_x\sum_y xP(x)yP(y) = \sum_x xP(X)\sum_y yP(Y) = E[X]E[Y]$$
\end{itemize}

La \textbf{Varianza} rappresenta la dispersione dei valori: $$Var(X)=E[(X-E[X])^2]$$  
Analogamente: 
$$ 
E[(X-E[X])^2]= E(X^2+E[X]^2-2E[X])=E[X^2]+E[X]^2-2E[X]E[X]=
E[X^2]-E[X]^2
$$

\newpage
\subsubsection{Esempio N.2 - Lotterie}
Consideriamo due possibili esempi di lotteria. Un'estrazione nella lotteria $L_1$ può risultare in:
\begin{itemize}
	\item Una perdità di 100 euro con probabilità $P=\frac{1}{4}$
	\item Una perdità di 200 euro con probabilità $P=\frac{1}{4}$
	\item Una vincità di 250 euro con probabilità
	$P=\frac{1}{2}$
\end{itemize}
Mentre in una lotteria $L_2$ ogni estrazione risulta in una vincita di 50 euro.

Calcoliamo il \textbf{valore atteso} di entrambe le lotterie $L_1$ ed $L_2$.

$$
E[L_1] = \frac{1}{4} \cdot(-100) + \frac{1}{4} (-200) + \frac{1}{2} (250) = -\frac{100}{4} - \frac{200}{4} + \frac{250}{2} = -25 - 50 + 125 = 50
$$

$$
E[L_2] = 1 \cdot 50 = 50
$$

Secondo i valori attesi di entrambe le lotterie, non esiste alcun vantaggio nel giocare in una lotteria o nell'altra. Adesso calcoliamo la \textbf{varianza} di entrambe le lotterie $L_1$ ed $L_2$
\\
$$
var(L_1) = E(L_1^2) - E(L_1)^2	
$$

Calcoliamo $E(L_1^2)$:
$$
E(L_1^2) = \frac{1}{4} \cdot(100)^2 + \frac{1}{4}\cdot(200)^2 + \frac{1}{2} \cdot(250)^2 = 2500 + 10000 + 31250 = 43750
$$

La varianza di $L_1$ risulta:
$$
var(L_1) = 43750 - 2500 = 41250
$$

Calcolare $var(L_2)$ è più semplice, in quanto sia $E(L_2^2)$ che $E(L_2)^2$ sono uguali a $50^2 = 2500$, quindi:
$$
var(L_2) = E(L_2^2) - E(L_2)^2 = 2500 - 2500 = 0
$$
Quest'ultimo risultato è scontato, in quanto per $L_2$ era possibile un solo caso con probabilità massima.

\subsubsection{Esercizio N.3 - Penitenziario}
In un penitenziario dove si trovano 3 prigionieri è annunciato che 2 di loro verranno liberati, senza tuttavia dire chi.
Uno dei prigionieri considera la possibilità di chiedere ad un amico secondino l’identità di uno dei prigionieri che sarà liberato (e che non sia lui stesso).\\ Il prigioniero, tuttavia, esita a porre tale domanda frenato dal seguente ragionamento.\\
Allo stato attuale la sua probabilità di essere liberato è $\frac{2}{3}$. Chiedendo al secondino (e ottenendo una risposta) la probabilità di essere rilasciato diventerebbe $\frac{1}{2}$, in quanto vi sarebbero due soli prigionieri
(uno dei quali è lui stesso) il cui destino è incerto e solo uno di essi verrebbe liberato. \\
Cosa c’è di sbagliato in questo ragionamento?\\

\begin{minipage}{0.50\textwidth}
	Definiamo innanzitutto la notazione:
	\begin{itemize}
		\item $Pr$ - Probabilità
		\item $P_i$ - Il prigioniero $i$ venga liberato
		\item $G_j$ - Il secondino dice che verrà liberato $P_j$
	\end{itemize}
\end{minipage}
\qquad
\begin{minipage}{0.45\textwidth}
	e le probabilità \textbf{note}:
	\begin{enumerate}
		\item $Pr[P_1 \wedge P_2] = \frac{1}{3}$
		\item $Pr[P_1 \wedge P_3] = \frac{1}{3}$
		\item $Pr[P_2 \wedge P_3] = \frac{1}{3}$
	\end{enumerate}
\end{minipage}

Supponiamo che il prigioniero che fa questo ragionamento è il prigioniero 1 e che il secondino gli riveli che verrà liberato il prigioniero 2, indicato con $G_2$. Le probabilità che ora abbiamo sono:
\begin{enumerate}
	\item $Pr[G_2|P_1 \wedge P_2] = 1$
	\item $Pr[G_2|P_1 \wedge P_3] = 0$
	\item $Pr[G_2|P_2 \wedge P_3] = \frac{1}{2}$
\end{enumerate}

Il prigioniero 1 quindi deve calcolare qual è la probabilità di essere liberato sapendo che $G_2$ sarà liberato, ovvero:
$$
Pr[P_1|G_2] = \frac{Pr[P_1 \wedge G_2]}{Pr[G_2]}
$$

\newpage
Calcoliamo $Pr[P_1 \wedge G_2]$. Esso può essere scritto come:
$$
Pr[P_1 \wedge G_2] = Pr[G_2 \wedge P_1] = Pr[G_2 \wedge P_1 \wedge P_2] + Pr[G_2 \wedge P_1 \wedge P_3] =
$$
$$
Pr[G_2|P_1 \wedge P_2] \cdot Pr[P_1 \wedge P_2] + Pr[G_2|P_1 \wedge P_3] \cdot Pr[P_1 \wedge P_3] =
$$
$$
1 \cdot Pr[P_1 \wedge P_2] + 0 \cdot Pr[P_1 \wedge P_3] = 1 \cdot \frac{1}{3} \quad\Rightarrow\quad Pr[P_1 \wedge G_2] = \frac{1}{3}
$$

Adesso calcoliamo $Pr[G_2]$, che può essere scritto come:
$$
Pr[G_2] = Pr[G_2 \wedge (P_1 \wedge P_2)] + Pr[G_2 \wedge (P_1 \wedge P_3)] + Pr[G_2 \wedge (P_2 \wedge P_3)] =
$$
$$
Pr[G_2|P_1 \wedge P_2] \cdot Pr[P_1 \wedge P_2] + Pr[G_2|P_1 \wedge P_3] \cdot Pr[P_1 \wedge P_3] + Pr[G_2|P_2 \wedge P_3] \cdot Pr[P_2 \wedge P_3] =
$$
$$
1 \cdot \frac{1}{3} + 0 \cdot \frac{1}{3} + \frac{1}{2} \cdot \frac{1}{3} = \frac{1}{3} + \frac{1}{6} = \frac{3}{6} = \frac{1}{2} \quad \Rightarrow \quad Pr[G_2] = \frac{1}{2}
$$

Quindi la probabilità che $P_1$ venga rilasciato sapendo che $G_2$ verrà rilasciato è:
$$
Pr[P_1|G_2] = \frac{Pr[P_1 \wedge G_2]}{Pr[G_2]} = \frac{\frac{1}{3}}{\frac{1}{2}} = \frac{2}{3}
$$

\subsubsection{Esempio N.4 - Dado non truccato}
Si supponga di lanciare un dado (non truccato) una volta al secondo, annotando quando esce il 6. Qual è il numero medio di lanci tra un 6 e il successivo? 

Possiamo risolvere questo problema con due metodi. Il primo metodo consiste nel definire la variabile aleatoria $t$, la cui probabilità è:
$$
Pr[t] = (\frac{5}{6})^{t-1}	\cdot \frac{1}{6}
$$
Il valore atteso di $t$ è:
$$
\sum_{t=1}^{\infty}t \cdot Pr[t] = \sum_{t=1}^{\infty}t \cdot (\frac{5}{6})^{t-1}	\cdot \frac{1}{6} = \sum_{t=1}^{\infty}[(t-1) (\frac{5}{6})^{t-1} \cdot \frac{1}{6} + (\frac{5}{6})^{t-1}	\cdot \frac{1}{6}] =
$$
$$
= \sum_{t=1}^{\infty}\ (t-1)\ (\frac{5}{6})^{t-1} \cdot \frac{1}{6} + \frac{1}{6}\  \sum_{t=1}^{\infty}\ (\frac{5}{6})^{t-1} = \sum_{t=1}^{\infty}\ (t-1)\ (\frac{5}{6})^{t-1} \cdot \frac{1}{6} + \frac{1}{6}\  \sum_{i=0}^{\infty}\ (\frac{5}{6})^i =
$$
$$
= \sum_{t=1}^{\infty}\ (t-1)\ (\frac{5}{6})^{t-1} \cdot \frac{1}{6} + \frac{1}{6} \cdot \frac{1}{1 - \frac{5}{6}} =
1 + \sum_{t=1}^{\infty}\ (t-1)\ (\frac{5}{6})^{t-1} \cdot \frac{1}{6} = 
$$
$$
= 1 + \frac{5}{6}\ \sum_{t=1}^{\infty}\ (t-1)\ (\frac{5}{6})^{t-2} \cdot \frac{1}{6} = 1 + \frac{5}{6}\  \sum_{t=1}^{\infty}\ t(\frac{5}{6})^{t-1} \cdot \frac{1}{6} =
$$
$$
= 1 + \frac{5}{6} E(t) \quad \Rightarrow \quad E(t) = 1 + \frac{5}{6} E(t)\quad \Rightarrow \quad \frac{1}{6} E(t) = 1
$$
$$
E(t) = 6
$$


Un metodo alternativo di svolgere questo esercizio è quello di usare un metodo "ricorsivo":

$$
E(t) = \frac{5}{6}(1 + E(t)) + \frac{1}{6} \cdot 1 = \frac{5}{6} + \frac{5}{6} E(t) + \frac{1}{6} = 1 + \frac{5}{6} E(t)
$$
$$
E(t) = 1 + \frac{5}{6} E(t)\quad \Rightarrow \quad \frac{1}{6} E(t) = 1 
$$
$$
E[t] = 6
$$
